\section{Discussion and limitations}\label{sec:disc}

\paragraph{What the AS sets measure.}
An AS counts as reached when it appears on an anomalous path that a peer exports to a collector.
Routes that no peer exports stay invisible.
Five of the ten collectors are at European IXPs and one at IX.br, so routes seen only at other IXPs are missing.
A path observation does not show the selected route or the forwarding state of its ASes, and a missing AS does not prove an outage.

\paragraph{Labels come from identifiers.}
The onset $t_0$ is the first observation at a collector, not the true start of an event.
The leak identifier misses leaks whose path triple was already seen in the same month, and it inherits the errors of inferred AS relationships~\cite{gao2001relationships,luckie2013relationships,jin2019problink}.
An adjacency that CAIDA does not list gets no relationship and is missed, and an outdated relationship type can create a false leak.

Novelty restarts every calendar month, so a recurring triple looks new on the first day of each month.
In 2025, the identifier date of \num{1174} of \num{8412} leaks is the first day of a month, against about 277 expected.
Of the \num{1346} leaks with 200 or more ASes, 278 have such a date, and 106 of the \num{1116} after the onset rule of \Cref{sec:prop}.
The onset rule removes 172 of these 278 leaks but 58 of the \num{1068} others, and R1 and R2 neither apply nor test it.
Without such leaks, C2 moves toward SAGE with history in a descriptive revision analysis (\Cref{app:protocol}).
% round 6 (layout, page limit): pointer to the paragraph "Month boundary", moved from sections/08-forecast.tex to sections/app-protocol.tex;
%   month_boundary.json C2 all -> no_day1: R1 +0.0054 -> +0.0066 (30 s), +0.0022 -> +0.0037 (60 s); R2 +0.0003 -> +0.0013, -0.0042 -> -0.0012
%   (both periods, both horizons). It removes first-of-month leaks, not the events that the onset rule removes, so it is not a test of that rule.
% src (round 6, panel-r5 NEW-13 / panel-r6 S8; review-r6 W5 item 5, question 6): arithmetic on the counts of the src comment
%   below (2025_leak/events.parquet): first-of-month 278 -> 106 kept, 172 removed = 61.9 %; other dates 1,346 - 278 = 1,068 -> 1,116 - 106 = 1,010
%   kept, 58 removed = 5.4 %. R1/R2: src/bgbench/forecast_protocol.py records observation_lag_s in the manifests without a filter; only the slow
%   regime filters observation_lag_s >= -60 (src/bgbench/baselines_eventset.py, curve_eventset.py). The audit sample shows such events directly:
%   leak_audit.json proportions.volume_day1.descriptors.route_visible_from_lookback_start 20 of 30 (radar_day1 0, radar_other 1, volume_other 2).
% src: paper/generated/onset_share_r1.json r1_cohort: observation_lag_s < -60 in 199 of the 3,543 R1 test events (manifest split == test joined to 2025_leak/events.parquet, exact cohort); R2 not computed because it needs the custodian event table, which is not read.
The month restart may favour repeat leakers, on which native history gains more by point estimate in a post hoc split (\Cref{sec:forecast-r1}).
% round 5 (panel-r5 NEW-4(c), NEW-16 / X2; domain B4, B7): "by point estimate in a post hoc split": the repeat-minus-first-seen native
%   difference includes zero under leaker x day ([-0.0309, +0.0009]) and in R2. The pointer now goes to the onset rule (sec:prop): the 1,116
%   events are 693 training and 423 test events, while sec:forecast-legacy prints 692 training events at 30 s (one has no 30-s input).
% round 4 (T6, A2): "may favour" stays. paper/generated/leaker_split.json measures where the gains sit (native@30 R1: first-seen -0.0112,
%   repeat -0.0260; class difference day blocks [-0.0236, -0.0062], leaker x day [-0.0309, +0.0009]), but no generated file crosses the
%   leaker class with the first-of-month identifier date, so the month effect on repeat leakers itself is not measured.
%   "recurring leakers" -> "repeat leakers" (one term, sections/08-forecast.tex "Repeat leakers").

We did not measure the precision of the identifier.
A revision audit of 120 R1 test leaks rates 95.7\,\% of the cohort relationship-consistent (\Cref{tab:leak-audit}).
This class re-runs the relationship test of the identifier one month later, so it does not measure precision (\Cref{app:protocol}).
% round 7: the sibling count (96 of 97) moved to app-protocol "Leak audit"; the independent-evidence sentence below replaces it here.
Beyond the identifier's own test, a Radar label with the same ASes or a short new route supports 71 of the 120 audited leaks.
% src (round 7, review-r6 W1 item 1): leak_audit.json independent.sample_totals.supported 71 (I3 50, I5 21); rules in app-protocol "Leak audit".
% src (round 7, review-r6 W1): paper/generated/leak_audit.json siblings.sample_totals.credible.independent 96, .sibling_only 1;
%   snapshots 20250901..20251201 (siblings.snapshots). Limits (AS2Org misses AS3356~AS209; intent not shown) in app-protocol "Leak audit".
Of the 60 audited volume-selected leaks, 10 have a changed or missing relationship and 11 are ambiguous.
R1 and R2 therefore describe the leaks that the identifier finds, matched to Radar or selected by volume.
% round 6 (panel-r6 NEW-2, NEW-9, W14 MADE_WORSE; review-r6 W1, item 1): the precision disclaimer of 2d0f7ef4 is restored; the class
%   "credible" is renamed "relationship-consistent" (round 7: the generated header and rows also use the new names);
%   scripts/paper_leak_audit.py gather() requires the event-month relationships to equal the identifier's classification, so the class re-runs the
%   identifier's own test one month later. "bounded" (a third sense of bound) is dropped. For the page limit, the weighted Wilson limits (82--99 %)
%   stay in tab:leak-audit, the 10-event sensitivity is its row "Without sampled likely FP", and the examples and the Radar-only strata are in
%   sections/app-protocol.tex "Leak audit".
% src (round 6): leak_audit.json cohort_estimate.weighted_by_stratum_size.credible 0.9571; proportions volume_day1 likely_false_positive 2 /
%   ambiguous 11, volume_other likely_false_positive 8 / ambiguous 0 (10 likely false and 11 ambiguous of the 60 audited volume-selected leaks).
% src (round 5, review-r5 W1, item 1; revision analysis, confirmatory false): paper/generated/leak_audit.json (scripts/paper_leak_audit.py,
%   first commit b644cf6b 2026-10-31 17:47 UTC; sample paper/generated/leak_audit_sample.csv, 120 rows). strata population/sampled: radar_day1
%   476/30, radar_other 2,968/30, volume_day1 65/30, volume_other 34/30 (cohort 3,543). cohort_estimate.weighted_by_stratum_size: credible 0.9571
%   (wilson_weighted_bounds [0.8225, 0.9879]), likely_false_positive 0.0038 [0.0017, 0.1185], ambiguous 0.0391 [0.0098, 0.1731] (the stratum
%   Wilson bounds with the same weights, "conservative descriptive bounds, not a joint 95 % interval"). proportions: radar_day1 and radar_other
%   29 credible + 1 ambiguous (A1_radar_conflict); volume_day1 17 credible, 2 likely FP, 11 ambiguous (A2_preexisting_path: route visible from
%   the look-back start, triples already new in the previous month); volume_other 22 credible, 8 likely FP (F1 5, F2 3).
%   sensitivity.C2@30 all +0.005436 -> excluding_sampled_likely_fp +0.005467; native@30 -0.023669 -> -0.023380 (both |change| < 0.001; published
%   values reproduced exactly). Round 6: the premise "We did not measure the precision of the identifier" is printed again (the audit does not
%   measure precision against ground truth), and the selected-candidate reading opens the next paragraph.
% Paragraph split (round 5): the audit and the catalogue checks form two paragraphs of at most six sentences.

For 83.7\,\% of the matched 2024 leaks, our reconstructed prefix count equals that of Radar.
The catalogues have known limits~\cite{grip,cloudflare2023hijack,liu2026}, so we use them to validate identifiers, not as ground truth.
Leaks that last longer than ten minutes are 3\,\% of 2024 events and carry thousands of prefixes.
Radar rarely labels them, so some may be policy changes, not leaks.
% src: final analysis/nongraph/2024/097_radar_prefix_agreement (971 of 1,160); propagation-facts.json leak_long_identifier
% src (round 4, panel-r4 NEW-9 / X2; author definition of 2026-10-31): data/corrected-release-20261016/eventset/2025_leak/events.parquet
%   (sha256 72e8b135...7163, read-only on 2026-10-31): 8,412 usable leaks (error null); asns_final_2h >= 200: 1,346, of which 278 have identifier date
%   event_id[6:8] == '01'; of these, observation_lag_s >= -60 (observed onset at most 60 s before the identifier, the onset rule of sec:prop):
%   1,116 events (693 train + 423 test splits), 106 with a first-of-month identifier date. 1,116 = forecast-facts.json
%   nongraph/2025/062_prefix_floor_boundary_v2 prefix_floor_summary.json large_leaks_ge200.events (test 423). The round-4 reviewer's 278 of 1,346 is
%   the population without the onset rule; both are printed now. Expected at uniform dates: 1,346 x 12/365 = 44 and 1,116 x 12/365 = 37.
%   "start on" -> identifier date; "This favours" -> "This may favour" (round-4 T6: no recurrence analysis yet).

\paragraph{Onset rules decide class comparisons.}
For valid and unknown origins, the first observation often marks an earlier, partial presence of the origin (\Cref{sec:prop}).
A comparison of early visibility between classes therefore needs an onset rule that neither differs by class nor selects by class.
% layout (round 5, page limit): "The first observation of a new origin is a natural onset, but its meaning differs by class" is folded into
%   the next sentence; sections/05-propagation.tex "Where the onset falls" gives the lags.
Our near-onset rule keeps different shares of each class, and a per-session clock as in \Cref{sec:recovery} exists only for the matched pairs.
We therefore report a sensitivity to the event definition, not a class difference that the onset explains.
In R1, 5.6\,\% of test events were first observed over 60\,s before the identifier.
For R2 this share needs the custodian table, so we did not compute it.
% src: sections/05-propagation.tex "Events with the onset near the change" (retention 84--94 % invalid, 28--46 % valid; fixed horizons only);
%   docs/paper/camera-ready-20261030/reviews-20261030/RESPONSE-ROUND1.md (re-anchoring t0 for all events not possible); round-2 review W1.

\paragraph{Selection shapes the RPKI classes.}
The label-free stratum removes the catalogue selection from invalid events only (\Cref{sec:prop}).
No similar stratum exists for valid and unknown origins, so we cannot remove it from those classes.
% layout (round 5, page limit): "Catalogues label invalid events far more often than valid ones, and the label-free stratum shows that this
%   selection raises the persistence floor of invalid events" is shortened; sections/05-propagation.tex "The catalogue shapes the invalid class"
%   gives the labelled shares (96--99 % against 10--36 %) and the floors.
% the frontier and RoVista-share caveats stay in sections 4 and 5 (consistency check, not evidence that a given AS filters); removed here for the page limit.

\paragraph{What the forecasts learn.}
The target sets hold ASes on paths to collector peers, so they can reflect these peers and their upstreams more than the wider topology.
A history of past events captures these recurring ASes, which can explain why it is hard to beat.
In a post hoc R1 analysis on reconstructed inputs, the same-leaker-first ranking adds 0.073 AP to it (\Cref{sec:forecast-r1}).
We did not separate forecasts of recurring ASes from forecasts of other ASes (\Cref{app:protocol}).
% src: paper/generated/leaker_split.json R1.same_leaker_baseline["30"].released_inputs.contrasts.same_leaker_then_history_minus_online_history_native
%   .ap.all 0.072939 (sections/08-forecast.tex "Repeat leakers"; post hoc, R1 only).

Every forecast uses an oracle trigger, the observed onset, so we do not show that a deployed detector could issue the forecast in time.
At 30\,s, 27 to 30\,\% of the first ten ranked ASes join (\Cref{tab:set-errors}).
Such a list could help an operator warn these ASes first, but we did not test this use.
% src (round 7, criterion M23(c)): set_errors_r2.tex P@10 at 30 s: R1 0.270-0.271, R2 0.271-0.295 (SAGE with history, online history
%   mean/median/native).
This works only if a detector fires within seconds of the onset.
% layout (round 5, page limit): "For the top ten, the native candidates of history add nothing" left; sections/08-forecast.tex "Secondary metrics" keeps it.
% src (round 4, NEW-5): paper/generated/set_errors_r2.json p_at_10_native_contrasts online_history_native_minus_online_history at 30 s:
%   R1 -0.0002 [-0.0013, +0.0011], R2 -0.0041 [-0.0070, -0.0012] (sections/08-forecast.tex "Secondary metrics", sections/app-protocol.tex).
% removed for the page limit (round-2 revision): a speculative sentence on graph models with a true-onset trigger or more collectors.

\paragraph{Event classes not covered.}
Link failures need last-seen times per edge, which our retained tables do not hold.
Our identifier compares the origins of the same prefix, so it misses a new more-specific prefix that no peer announced before.
% src (round 5, review-r5 W4 / item 4, panel-r5 W6(a)): src/bgbench/prefix_events.py origin_changes (an announcement whose origin
%   differs from the previous origin of the same peer and the same prefix; prev_origin = shift over (peer_ip, prefix)). An already observed
%   more-specific prefix can still have an origin change, so sub-prefix events are not all excluded.
Forged-origin hijacks keep the legitimate origin and have no label source, and DFOH~\cite{holterbach2024dfoh} would be the external check.
% DFOH detects forged-origin hijacks (refs.bib title "A System to Detect Forged-Origin BGP Hijacks"); round-3 review NEW-16 / R2 C10.
